\section{Conclusions and perspectives}
We have presentd our vision of a decentralized architecture for virtual environments. Thanks to a N-dimensional Vorono� based p2p network, called RayNet, we can distibute the exchange of data and the computation of heavy tasks over the different nodes present in the virtual space. Descriptors presented in this papers allows us first managing with respect to view-dependency criteria the adaptive streaming of 3D models from neighborhood peers with the optimal resolution, and secondly filtering the computation of tasks such as detection collision according to the proximity between peers. Moreover, we have presented a management policy for objects to be able to use them, to control them using a user-defined software components, or to drop them in the virtual space. Finally, we describe the functionnalities of our navigator, allowing the visualization of Web 2.0 components into the virtual world , the embedding of the navigator into web pages thanks to interactive texuring, the modeling of 3D contents integrated to the viewer, and the possibility to interact with the universe from clients with low ressources such as mobile phones.\\
The use of a N-Dimensional p2p network allow us to extend the connectivity not only according to the proximity of peers in the virtual environment, but according to social or semantic proximity between object or avatar nodes. As the metaverse belong to all users, this project is fully opensource, and we hope that community will help us to improve our architecture and evaluate in full-size its effectiveness.  
       